\documentclass[11pt]{article}
\usepackage[margin=1in]{geometry}
\usepackage{amsmath,amssymb,amsthm}
\usepackage{booktabs}
\usepackage{array}
\usepackage[T1]{fontenc}
\usepackage{lmodern}
\usepackage{microtype}

\newcommand{\XiN}{\mathsf{Xi}}
\newcommand{\Xiw}{\mathsf{Xi}_{W}}
\newcommand{\Xig}{\mathsf{Xi}_{G}}
\newcommand{\Xic}{\mathsf{Xi}_{C}}
\newcommand{\Rtwo}{R^2}
\newcommand{\Cp}{C_{\perp}}
\newcommand{\Wp}{W_{\perp}}
\newcommand{\Gp}{G_{\perp}}
\newtheorem{lemma}{Lemma}
\newtheorem{theorem}{Theorem}

\title{Phase--Radius Decoupling in a Finite V59 Residual}
\author{Liang Wang\\
School of Mathematics and Statistics\\
Huazhong University of Science and Technology (HUST), Wuhan, China}
\date{27 August 2026}

\begin{document}
\maketitle

\begin{abstract}
TPC-270 found a strong finite variation in an endpoint-normalized residual
radius, but did not record its signed scalar phase and its two norm lanes in
the same certificate.  We introduce
\(\Cp=\langle(I-P_3)w,(I-P_3)g\rangle\), \(\Wp=\|(I-P_3)w\|^2\), and
\(\Gp=\|(I-P_3)g\|^2\), together with rational sixth-power coordinates
\(\XiN=(\Wp\Gp)^3/N^{10}\), \(\Xiw=\Wp^3/N^5\),
\(\Xig=\Gp^3/N^5\), and \(\Xic=|\Cp|^6/N^{10}\).  The exact identities
\(\XiN=\Xiw\Xig\) and \(\XiN/\Xic=|\kappa|^{-6}\) separate radius size,
lane attribution, and phase alignment.  On the registered TPC-269 finite
interface, all six base rows and three matched profile controls have
negative-real-axis scalar phase.  Nevertheless the dyadic radius ratios have
pattern DROP--RISE--RISE--DROP; the 96 to 192 ratio exceeds 23 while the
source lane falls below 1/8 and the output lane rises above 230.  This is a
numerically certified finite phase--radius decoupling audit, not an asymptotic
phase theorem, radius estimate, arithmetic $L^2$ result, or twin-prime proof.
\end{abstract}

\section{Position and claim firewall}

The TPC chain studies a literal finite V59 physical operator containing the
prime shell, outer $q$ weight, unit masks, deleted diagonal, source beta,
registered cutoff $z_N=\lfloor\log N\rfloor$, and a rank-three four-block
projection.  TPC-269 added a convex path between two kernel profiles.  TPC-270
then normalized the positive residual radius by the endpoint exponent $5/3$
and found a finite DROP--RISE--RISE--DROP pattern.

The status of the present paper is
\[
\texttt{NUMERICALLY\_CERTIFIED\_FINITE\_PHASE\_RADIUS\_DECOUPLING\_AUDIT}.
\]
The word ``finite'' is essential: the registered rows are a diagnostic
interface, not a proved asymptotic sequence.  No finite ratio is promoted to a
power saving or a source-level counterexample.

\section{Projected residual coordinates}

Let $P_3$ be the orthogonal projection onto the three declared block
contrasts.  On each finite row put
\[
 \Cp=\langle(I-P_3)w_N,(I-P_3)g_{N,\theta}\rangle,\qquad
 \Wp=\|(I-P_3)w_N\|^2,\qquad
 \Gp=\|(I-P_3)g_{N,\theta}\|^2.
\]
The residual radius product and normalized signed coordinate are
\[
 \Rtwo=\Wp\Gp,\qquad
 \kappa=\frac{\Cp}{\sqrt{\Wp\Gp}}.
\]
The producer certifies $\Wp>0$ and $\Gp>0$ on all nine rows.  It also
certifies that every displayed $\Cp$ interval lies strictly below zero.

\section{Exact lane factorization}

The algebraic number $N^{5/3}$ is unnecessary at finite scale.  Define
\[
 \XiN=\frac{(\Rtwo)^3}{N^{10}},\qquad
 \Xiw=\frac{\Wp^3}{N^5},\qquad
 \Xig=\frac{\Gp^3}{N^5},\qquad
 \Xic=\frac{|\Cp|^6}{N^{10}}.
\]

\begin{lemma}[finite identities]
On every row with $\Wp,\Gp,\Cp\ne0$,
\[
 \XiN=\Xiw\Xig,
 \qquad
 \frac{\XiN}{\Xic}
 =\frac{(\Wp\Gp)^3}{|\Cp|^6}=|\kappa|^{-6}.
\]
\end{lemma}

\noindent
Indeed, the first equality is direct multiplication and the second follows
from $(\sqrt{\Wp\Gp}/|\Cp|)^6=|\kappa|^{-6}$.  Thus a radius change can be
attributed to the source lane $\Xiw$ and output lane $\Xig$, while the ratio
to $\Xic$ measures phase misalignment.  For positive outward intervals,
cubing is monotone and division uses the lower numerator over upper
denominator and vice versa.  The certificate stores these outward images.

\section{Certified finite result}

The base registry is
\[
 (N,H,Q)=(64,15,4),(96,20,5),(128,24,5),
 (192,32,6),(256,38,6),(384,50,7).
\]
The cutoff is $z_N=\lfloor\log N\rfloor$ on these rows.  The four dyadic
lane records are shown below; all endpoints are outward decimal enclosures.

\begin{table}[h]
\centering
\scriptsize
\begin{tabular}{c@{\quad}c@{\quad}c@{\quad}c}
\toprule
pair & $\Xiw$ ratio & $\Xig$ ratio & $\XiN$ ratio\\
\midrule
64$\to$128 & [0.365585,0.365733] & [0.633926,0.633926] & [0.231754,0.231847]\\
96$\to$192 & [0.103825,0.103864] & [230.7698,230.7699] & [23.9597,23.9686]\\
128$\to$256 & [0.466740,0.466928] & [15.3653,15.3654] & [7.17162,7.17449]\\
192$\to$384 & [1.10084,1.10125] & [0.729362,0.729363] & [0.802913,0.803208]\\
\bottomrule
\end{tabular}
\caption{Source, output, and endpoint-normalized radius lane ratios.}
\end{table}

\begin{theorem}[finite phase--radius decoupling]
On the six base rows, all scalar intervals $\Cp$ are strictly negative, so the
phase label is negative real throughout.  Simultaneously, the four dyadic
radius ratios have classifications
\[
 \text{DROP}_{<1/4},\quad \text{RISE}_{>23},\quad
 \text{RISE}_{>7},\quad \text{DROP}_{(3/4,1)}.
\]
The $96\to192$ rise is output-lane dominated: its source-lane ratio is below
$1/8$ and its output-lane ratio is above $230$.
\end{theorem}

\noindent
The classifications are separated from their thresholds by the stored
intervals.  In particular, the largest radius spike occurs while the signed
phase label is unchanged.  This is an attribution statement about the finite
registry, not a claim that phase and radius are statistically independent.

At $N=96,128,256$, the matched $\theta=1/2$ controls preserve the negative
phase and leave the source lane unchanged because $w_N$ is unchanged.  Their
output-lane ratios are respectively below $0.9$, while their radius ratios
are enclosed in $(1/2,3/4)$.  This provides a profile control for the lane
interpretation.

\section{Interpretation and limits}

The second identity in Lemma 1 explains why a small signed scalar can coexist
with a large radius envelope: the phase coordinate $|\kappa|$ can be small,
and the amplification factor is exactly $|\kappa|^{-6}$.  The lane identity
then separates two mechanisms that a radius-only table conflates.  At
$96\to192$, the output norm lane supplies the large factor even though the
source-normalized lane decreases.  At $64\to128$, both normalized lanes
decrease and the radius drops below $1/4$ without a phase sign change.

Four promotions are explicitly disallowed.  The six rows do not establish a
uniform source-level sequence; the phase sign does not establish an eventual
phase sector; the radius ratios do not prove or refute a bound
$R_N\ll N^{5/3-\delta}$; and the radius product is not an arithmetic $L^2$
estimate or a signed four-packet reassembly.  Accordingly the fixed-power
credit is zero and full Gate B remains open.

\section{Conclusion}

TPC-271 adds a joint coordinate system to the TPC-267--270 finite chain.  The
exact lane factorization and the independently audited records show a clean
finite phase--radius decoupling: negative-real scalar phase persists while the
normalized radius rises by more than 23, driven by the output residual lane.
The next source-level task is a uniform signed-phase estimate coupled to an
explicit radius-lane bound.  Until that estimate is proved, no asymptotic or
twin-prime conclusion follows.

\begin{thebibliography}{9}
\bibitem{HW}
G. H. Hardy and E. M. Wright, \emph{An Introduction to the Theory of Numbers},
6th ed., Oxford University Press, 2008.
\bibitem{MV}
H. L. Montgomery and R. C. Vaughan, \emph{Multiplicative Number Theory I},
Cambridge University Press, 2007.
\bibitem{TPC270}
L. Wang, ``Cross-scale endpoint-normalized radius in a finite V59 residual,''
TPC-270 project release, 2026.
\end{thebibliography}

\end{document}
